
\documentclass[11pt,a4paper]{article}
\usepackage[T1]{fontenc}
\usepackage{lmodern}
\usepackage{amsmath,amssymb,amsfonts,mathtools}
\usepackage{graphicx,booktabs,longtable,array,tabularx,adjustbox}
\usepackage{geometry,microtype,fancyhdr,lastpage,etoolbox,enumitem}
\usepackage{newunicodechar,textcomp}
\usepackage{hyperref}
\hypersetup{hidelinks,pdfcreator={XeLaTeX / QM1 academic edition}}
\geometry{left=20mm,right=20mm,top=22mm,bottom=23mm}
\setcounter{secnumdepth}{2}
\setlength{\parindent}{0pt}
\setlength{\parskip}{4pt plus 1pt}
\setlength{\emergencystretch}{2em}
\setlength{\tabcolsep}{3pt}
\renewcommand{\arraystretch}{1.16}
\setlist{nosep,topsep=3pt,itemsep=2pt,leftmargin=7mm}
\AtBeginEnvironment{longtable}{\footnotesize\setlength{\tabcolsep}{2.2pt}}
\newcommand{\tightlist}{\setlength{\itemsep}{0pt}\setlength{\parskip}{0pt}}
\providecommand{\Cov}{\operatorname{Cov}}
\providecommand{\Var}{\operatorname{Var}}
\providecommand{\Corr}{\operatorname{Corr}}
\providecommand{\MAD}{\operatorname{MAD}}
\providecommand{\E}{\mathbb{E}}
\providecommand{\R}{\mathbb{R}}
\newunicodechar{—}{\textemdash}
\newunicodechar{–}{\textendash}
\newunicodechar{€}{\texteuro}
\newunicodechar{≤}{\ensuremath{\leq}}
\newunicodechar{≥}{\ensuremath{\geq}}
\newunicodechar{×}{\ensuremath{\times}}
\newunicodechar{÷}{\ensuremath{\div}}
\newunicodechar{Σ}{\ensuremath{\Sigma}}
\newunicodechar{σ}{\ensuremath{\sigma}}
\newunicodechar{μ}{\ensuremath{\mu}}
\newunicodechar{Δ}{\ensuremath{\Delta}}
\newunicodechar{∑}{\ensuremath{\sum}}
\newunicodechar{≈}{\ensuremath{\approx}}
\newunicodechar{√}{\ensuremath{\surd}}
\newunicodechar{≠}{\ensuremath{\ne}}
\newunicodechar{→}{\ensuremath{\to}}
\newunicodechar{∞}{\ensuremath{\infty}}
\newunicodechar{∈}{\ensuremath{\in}}
\newunicodechar{−}{\ensuremath{-}}
\newunicodechar{²}{\ensuremath{^{2}}}
\newunicodechar{³}{\ensuremath{^{3}}}
\newunicodechar{₀}{\ensuremath{_{0}}}
\newunicodechar{₁}{\ensuremath{_{1}}}
\newunicodechar{₂}{\ensuremath{_{2}}}
\newunicodechar{₃}{\ensuremath{_{3}}}
\newunicodechar{₄}{\ensuremath{_{4}}}
\newunicodechar{ₜ}{\ensuremath{_{t}}}
\newunicodechar{ₚ}{\ensuremath{_{p}}}
\newunicodechar{ₙ}{\ensuremath{_{n}}}
\newunicodechar{ₓ}{\ensuremath{_{x}}}
\newunicodechar{ᵢ}{\ensuremath{_{i}}}
\newunicodechar{ŷ}{\ensuremath{\hat y}}
\newunicodechar{Ŷ}{\ensuremath{\widehat Y}}
\newunicodechar{ₘ}{\ensuremath{_{m}}}
\newunicodechar{ₑ}{\ensuremath{_{e}}}
\newunicodechar{ₐ}{\ensuremath{_{a}}}
\newunicodechar{ₖ}{\ensuremath{_{k}}}
\newunicodechar{ₛ}{\ensuremath{_{s}}}
\newunicodechar{ₗ}{\ensuremath{_{l}}}
\newunicodechar{ₕ}{\ensuremath{_{h}}}
\newunicodechar{ₒ}{\ensuremath{_{o}}}
\newunicodechar{₊}{\ensuremath{_{+}}}
\newunicodechar{₋}{\ensuremath{_{-}}}
\newunicodechar{₌}{\ensuremath{_{=}}}
\newunicodechar{₍}{\ensuremath{_{(}}}
\newunicodechar{₎}{\ensuremath{_{)}}}
\newunicodechar{⁻}{\ensuremath{^{-}}}
\newunicodechar{⁺}{\ensuremath{^{+}}}
\newunicodechar{⁰}{\ensuremath{^{0}}}
\newunicodechar{¹}{\ensuremath{^{1}}}
\newunicodechar{⁴}{\ensuremath{^{4}}}
\newunicodechar{⁵}{\ensuremath{^{5}}}
\newunicodechar{⁶}{\ensuremath{^{6}}}
\newunicodechar{⁷}{\ensuremath{^{7}}}
\newunicodechar{⁸}{\ensuremath{^{8}}}
\newunicodechar{⁹}{\ensuremath{^{9}}}
\makeatletter
\renewcommand{\section}{\@startsection{section}{1}{0pt}{2.0ex plus .7ex minus .3ex}{1.0ex}{\large\bfseries}}
\renewcommand{\subsection}{\@startsection{subsection}{2}{0pt}{1.8ex plus .5ex minus .2ex}{.7ex}{\normalsize\bfseries}}
\renewcommand{\subsubsection}{\@startsection{subsubsection}{3}{0pt}{1.3ex plus .4ex minus .2ex}{.5ex}{\normalsize\bfseries}}
\makeatother
\pagestyle{fancy}
\fancyhf{}
\renewcommand{\headrulewidth}{0.35pt}
\renewcommand{\footrulewidth}{0pt}
\fancyhead[L]{\small\textsc{KEDGE Business School --- QM1}}
\fancyhead[R]{\small\textsc{2026--2027}}
\fancyfoot[L]{\small Badr ABOUFARASSE, PhD}
\fancyfoot[R]{\small\thepage\ /\ \pageref{LastPage}}
\setlength{\headheight}{14pt}
\setkeys{Gin}{keepaspectratio}
\begin{document}

\begin{center}
{\small\textsc{KEDGE Business School \quad EBP1 \quad Quantitative Methods 1}}\\[6pt]
{\Large\bfseries Fiche 6 — Nombres-indices}\\[6pt]
{\normalsize\bfseries Énoncés des exercices}\\[6pt]
{\small Badr ABOUFARASSE, PhD \quad --- \quad 2026--2027}
\end{center}
\vspace{2mm}\hrule\vspace{6mm}

\hypertarget{applications-fondamentales}{%
\section{Applications fondamentales}\label{applications-fondamentales}}

Maîtriser les définitions et les calculs essentiels du cours.

\emph{Application du cours --- 20 min}

\hypertarget{exercice-1-indices-uxe9luxe9mentaires-rebasage-et-propriuxe9tuxe9s}{%
\subsection{Exercice 1 -- Indices élémentaires, rebasage et propriétés}\label{exercice-1-indices-uxe9luxe9mentaires-rebasage-et-propriuxe9tuxe9s}}

Un service coûte 80 euros en 2023, 92 euros en 2024 et 101.20 euros en 2025.

\begin{enumerate}
\item
  Calculer les indices 2024/2023, 2025/2024 et 2025/2023, base 100.
\item
  Vérifier la circularité.
\item
  Calculer l'indice inverse 2023/2025 et vérifier la réversibilité.
\item
  Rebaser la série pour prendre 2024 comme base 100.
\item
  Interpréter correctement l'indice 126.5.
\end{enumerate}

\emph{Application du cours --- 25 min}

\hypertarget{compluxe9ment-1-duxe9composition-compluxe8te-prixquantituxe9s}{%
\subsection{Complément 1 --- Décomposition complète prix--quantités}\label{compluxe9ment-1-duxe9composition-compluxe8te-prixquantituxe9s}}

Un service achats observe :

\begin{longtable}[]{@{}lllll@{}}
\toprule\noalign{}
Produit & p0 & q0 & p1 & q1 \\
\midrule\noalign{}
\endhead
\bottomrule\noalign{}
\endlastfoot
A & 4 & 100 & 5 & 90 \\
B & 10 & 50 & 11 & 60 \\
C & 20 & 20 & 18 & 30 \\
\end{longtable}

\begin{enumerate}
\tightlist
\item
  Calculez les quatre totaux Σp0q0, Σp1q0, Σp0q1 et Σp1q1.
\item
  Calculez les indices de valeur, de prix de Laspeyres, Paasche et Fisher.
\item
  Interprétez l'écart entre valeur et prix.
\end{enumerate}

\textbf{Indication méthodologique :}

Complétez les quatre colonnes de valeurs croisées avant tout rapport.

\hypertarget{pruxe9paration-aux-uxe9valuations}{%
\section{Préparation aux évaluations}\label{pruxe9paration-aux-uxe9valuations}}

Résoudre des exercices de contrôle continu et de partiel en autonomie.

\emph{Préparation CC / partiel --- 35 min}

\hypertarget{exercice-2-laspeyres-paasche-fisher-et-factor-reversal}{%
\subsection{Exercice 2 -- Laspeyres, Paasche, Fisher et factor reversal}\label{exercice-2-laspeyres-paasche-fisher-et-factor-reversal}}

\begin{longtable}[]{@{}lllll@{}}
\toprule\noalign{}
& {\(p_0\)} & {\(q_0\)} & {\(p_1\)} & {\(q_1\)} \\
\midrule\noalign{}
\endhead
\bottomrule\noalign{}
\endlastfoot
Produit X & 4 & 100 & 5 & 90 \\
Produit Y & 10 & 40 & 12 & 50 \\
\end{longtable}

\begin{enumerate}
\item
  Calculer les quatre agrégats {\(\sum p_0q_0\)}, {\(\sum p_1q_0\)}, {\(\sum p_0q_1\)}, {\(\sum p_1q_1\)}.
\item
  Calculer les indices de prix de Laspeyres, Paasche et Fisher.
\item
  Calculer les indices de quantité de Laspeyres, Paasche et Fisher.
\item
  Calculer l'indice de valeur.
\item
  Vérifier la propriété de factor reversal de Fisher sur une échelle base 100.
\end{enumerate}

\emph{Préparation CC / partiel --- 45 min}

\hypertarget{exercice-4-indices-complets-sur-quatre-produits-type-partiel}{%
\subsection{Exercice 4 --- Indices complets sur quatre produits --- type partiel}\label{exercice-4-indices-complets-sur-quatre-produits-type-partiel}}

Pour quatre produits, les prix et quantités sont observés à deux dates. Base 0 = 100.

\begin{longtable}[]{@{}lllll@{}}
\toprule\noalign{}
Produit / Product & p₀ & q₀ & p₁ & q₁ \\
\midrule\noalign{}
\endhead
\bottomrule\noalign{}
\endlastfoot
A & 53 & 20 & 62 & 25 \\
B & 41 & 28 & 52 & 34 \\
C & 75 & 35 & 82 & 38 \\
D & 25 & 30 & 38 & 33 \\
\end{longtable}

\begin{enumerate}
\tightlist
\item
  Construire les quatre colonnes de produits p₀q₀, p₁q₀, p₀q₁, p₁q₁ et leurs totaux.
\item
  Calculer les indices de prix et de quantité de Laspeyres, Paasche et Fisher.
\item
  Calculer l\textquotesingle indice de valeur et vérifier la propriété de Fisher. Interpréter les trois indices.
\end{enumerate}

\hypertarget{approfondissement-facultatif}{%
\section{Approfondissement facultatif}\label{approfondissement-facultatif}}

Aller plus loin dans l'interprétation et les applications professionnelles.

\emph{Approfondissement --- 15 min}

\hypertarget{exercice-3-indices-chauxeenuxe9s-et-duxe9flation}{%
\subsection{Exercice 3 -- Indices chaînés et déflation}\label{exercice-3-indices-chauxeenuxe9s-et-duxe9flation}}

Les indices annuels chaînés sont 105, 108 et 97.

\begin{enumerate}
\item
  Calculer l'indice cumulé sur les trois périodes.
\item
  Interpréter le résultat et expliquer pourquoi il ne faut pas additionner les variations annuelles.
\item
  Un chiffre d'affaires nominal vaut 2.4 millions d'euros lorsque l'indice des prix vaut 120. Calculer la valeur réelle en prix de la période de base.
\end{enumerate}

\emph{Approfondissement --- 25 min}

\hypertarget{compluxe9ment-2-changement-de-base-inflation-et-performance-ruxe9elle}{%
\subsection{Complément 2 --- Changement de base, inflation et performance réelle}\label{compluxe9ment-2-changement-de-base-inflation-et-performance-ruxe9elle}}

Le chiffre d'affaires nominal vaut 500 k€, 560 k€ et 630 k€. L'indice des prix correspondant (année 1 = 100) vaut 100, 108 et 120.

\begin{enumerate}
\tightlist
\item
  Calculez le chiffre d'affaires réel en euros de l'année 1.
\item
  Calculez les croissances nominale et réelle entre les années 1 et 3.
\item
  Rebasez l'indice des prix avec année 2 = 100.
\item
  Répondez au responsable qui affirme que l'activité a progressé de 26 \%.
\end{enumerate}

\textbf{Indication méthodologique :}

Déflatez en divisant la valeur nominale par indice/100. Le changement de base ne modifie pas la trajectoire d'inflation.

\end{document}
